% First NFCorpus recipe: pairs/SDPA, LoRA, .1 BCE + 1 pairwise.
% Same-name cross-row references preserve the graph; no detach.
% FP32 master parameters and scores/loss; supported compute uses BF16 autocast.
% Shared tensor geometry and palette. Cell counts/shades are illustrative.
\documentclass[tikz,border=2mm]{standalone}
\usepackage[UTF8,fontset=none]{ctex}
\setmainfont{DejaVu Sans}
\setsansfont{DejaVu Sans}
\setmonofont{DejaVu Sans Mono}
\setCJKmainfont{Noto Sans CJK SC}
\setCJKsansfont{Noto Sans CJK SC}
\setCJKmonofont{Noto Sans Mono CJK SC}
\usepackage{amsmath,amssymb}
\usepackage{tikz}
\usetikzlibrary{arrows.meta,calc,backgrounds}
\usepackage{hyperref}
\hypersetup{pdfauthor={shichao},pdfcreator={TeX}}
\definecolor{Ink}{HTML}{273648}
\definecolor{Muted}{HTML}{607086}
\definecolor{Data}{HTML}{408AA4}
\definecolor{Hidden}{HTML}{8C73B5}
\definecolor{Weight}{HTML}{D89548}
\definecolor{Score}{HTML}{C96E78}
\definecolor{Target}{HTML}{7E8B9B}
\definecolor{Chosen}{HTML}{388878}
\newcommand{\Small}{\fontsize{8.5}{10.5}\selectfont}
\newcommand{\Body}{\fontsize{10}{12}\selectfont}
\newcommand{\Heading}{\fontsize{11}{13}\selectfont\bfseries}
\tikzset{
  flow/.style={-{Latex[length=2mm,width=1.35mm]},draw=Ink,line width=.75pt},
  grad/.style={-{Latex[length=2mm,width=1.35mm]},draw=Score,line width=.85pt,dashed},
  block/.style={draw=Ink!25,fill=white,rounded corners=2pt,line width=.55pt,
                align=center,font=\Body,inner sep=5pt},
  note/.style={font=\Small,text=Muted,align=center},
  heading/.style={font=\Heading},
  dim/.style={font=\Small,text=Muted,align=center}}
\newcommand{\Band}[5]{%
  \begin{scope}[on background layer]
    \path[fill=#5!6,rounded corners=5pt] (#1,#2) rectangle (#3,#4);
  \end{scope}}
% name,x,y,width,height,rows,columns,color,symbol
\newcommand{\Tensor}[9]{%
  \coordinate (#1) at (#2,#3);
  \coordinate (#1-west) at ({#2-#4/2},#3);
  \coordinate (#1-east) at ({#2+#4/2},#3);
  \coordinate (#1-north) at (#2,{#3+#5/2});
  \coordinate (#1-south) at (#2,{#3-#5/2});
  \foreach \r in {1,...,#6}{\foreach \c in {1,...,#7}{%
    \pgfmathtruncatemacro{\tone}{28+mod(\r+2*\c,3)*22}
    \path[draw=#8!20,fill=#8!\tone,line width=.2pt,rounded corners=.6pt]
      ({#2-#4/2+(\c-1)*#4/#7+.015},{#3+#5/2-\r*#5/#6+.015})
      rectangle ({#2-#4/2+\c*#4/#7-.015},{#3+#5/2-(\r-1)*#5/#6-.015});}}
  \node[font=\Body] at (#2,{#3-#5/2-.31}) {$#9$};}
\newcommand{\Dim}[3]{\node[dim] at (#1,#2) {$#3$};}
\newcommand{\HiddenTensor}[3]{%
  \path[draw=Hidden!30,fill=white,rounded corners=1pt]
    (#2-.82,#3-.57) rectangle (#2+1.18,#3+.93);
  \Tensor{#1}{#2}{#3}{2}{1.5}{5}{4}{Hidden}{#1}}
\newcommand{\BackbonePanel}[2]{%
  \Band{.1}{#1-2.4}{27.9}{#1+3.3}{Hidden}
  \node[anchor=west,heading,text=Hidden!85!Ink] at (.4,#1+2.84)
    {1. 候选微批：输入、主干与末尾读出};
  \node[anchor=east,note] at (27.5,#1+2.84) {#2};
  \Tensor{X}{1.8}{#1}{1.5}{1.05}{3}{5}{Data}{X}
  \Dim{1.8}{#1-1.22}{[b,L]\;\mathrm{int64}}
  \node[note] at (1.8,#1+1.4) {每行一个候选序列};
  \node[block,minimum width=2cm,minimum height=1cm] (embed) at (5.25,#1)
    {Embedding};
  \HiddenTensor{E}{8.1}{#1}
  \Dim{8.1}{#1-1.48}{[b,L,D]}
  \node[note] at (8.1,#1+1.4) {叠层 $b$；面 $L\times D$};
  \node[block,minimum width=3.5cm,minimum height=1.8cm,font=\Small]
    (decoder) at (12.2,#1)
    {\textbf{Qwen3 Decoder $\times28$}\\[4pt]
     q/k/v/o LoRA\\[4pt]含最终 RMSNorm};
  \node[note] at (12.2,#1+1.93)
    {$\textcolor{Data}{M}:[b,L]$\quad $\textcolor{Target}{P}:[b,L]$\\[-1pt]
     $M$: padding 支持；$P$: 逻辑位置};
  \draw[flow] (12.2,#1+1.33) -- (decoder.north);
  \HiddenTensor{H}{16.5}{#1}
  \Dim{16.5}{#1-1.48}{[b,L,D]}
  \node[block,minimum width=2.9cm,minimum height=1.05cm,font=\Small]
    (gather) at (20.4,#1)
    {\textbf{Gather}\\[5pt]$H[j,\ell_j-1,:]$};
  \node[note] at (20.4,#1-1.35) {$\ell_j$：该行有效长度};
  \Tensor{Z}{25.2}{#1}{2}{1.05}{3}{4}{Hidden}{Z}
  \Dim{25.2}{#1-1.22}{[b,D]}
  \node[note] at (25.2,#1+1.4) {末尾有效后缀位置};
  \draw[flow] (X-east) -- (embed.west);
  \draw[flow] (embed.east) -- (E-west);
  \draw[flow] (9.28,#1) -- (decoder.west);
  \draw[flow] (decoder.east) -- (H-west);
  \draw[flow] (17.68,#1) -- (gather.west);
  \draw[flow] (gather.east) -- (Z-west);
  \node[note,anchor=west] at (.4,#1-2.12)
    {$M,P$ 为 int64；$M$ 不是完整因果矩阵。注意力：$Q:[b,16,L,128]$，$K,V:[b,8,L,128]$。};}
\newcommand{\HeadPanel}[1]{%
  \Band{.1}{#1-2.4}{27.9}{#1+2.6}{Weight}
  \node[anchor=west,heading,text=Weight!85!Ink] at (.4,#1+2.14)
    {2. 独立评分头：$D\to h\to1$};
  \node[anchor=east,note] at (27.5,#1+2.14)
    {左侧 $Z$ 承接上行；$s_{\mathrm{chunk}}$ 供下行合并微批};
  \Tensor{Zhead}{1.8}{#1}{2}{1.05}{3}{4}{Hidden}{Z}
  \Dim{1.8}{#1-1.22}{[b,D]}
  \node[font=\fontsize{15}{17}\selectfont] at (3.45,#1) {$\times$};
  \Tensor{Wone}{4.7}{#1}{1.25}{2}{4}{3}{Weight}{W_1^{\mathsf T}}
  \Dim{4.7}{#1-1.68}{[D,h]}
  \Tensor{U}{8.9}{#1}{1.25}{1.05}{3}{3}{Hidden}{U}
  \Dim{8.9}{#1-1.22}{[b,h]}
  \draw[flow] (Wone-east) -- (U-west);
  \node[draw=Weight!35,fill=Weight!12,rounded corners=1pt,font=\Small,
    inner sep=3pt] at (6.8,#1+.55) {$+b_1\;[h]$};
  \node[block,minimum width=1.6cm,minimum height=.8cm] (gelu) at (11.45,#1) {GELU};
  \draw[flow] (U-east) -- (gelu.west);
  \Tensor{A}{14}{#1}{1.25}{1.05}{3}{3}{Hidden}{A}
  \Dim{14}{#1-1.22}{[b,h]}
  \draw[flow] (gelu.east) -- (A-west);
  \node[font=\fontsize{15}{17}\selectfont] at (15.4,#1) {$\times$};
  \Tensor{Wtwo}{16.8}{#1}{.35}{1.25}{3}{1}{Weight}{W_2^{\mathsf T}}
  \Dim{16.8}{#1-1.32}{[h,1]}
  \Tensor{v}{20.3}{#1}{.35}{1.05}{3}{1}{Score}{v}
  \Dim{20.3}{#1-1.22}{[b,1]}
  \draw[flow] (Wtwo-east) -- (v-west);
  \node[draw=Weight!35,fill=Weight!12,rounded corners=1pt,font=\Small,
    inner sep=3pt] at (18.5,#1+.55) {$+b_2\;[1]$};
  \node[block,minimum width=2.2cm,minimum height=1.05cm,font=\Small]
    (squeeze) at (23.1,#1) {squeeze($-1$)\\[4pt]\texttt{.float()}};
  \draw[flow] (v-east) -- (squeeze.west);
  \Tensor{chunk}{26.25}{#1}{.35}{1.05}{3}{1}{Score}{\boldsymbol{s}_{\mathrm{chunk}}}
  \Dim{26.25}{#1-1.22}{[b]\;\mathrm{FP32}}
  \draw[flow] (squeeze.east) -- (chunk-west);
  \node[note,anchor=west] at (.4,#1-2.08)
    {实际存储：$W_1:[h,D]$、$W_2:[1,h]$；图中矩乘展示转置。$b_1,b_2$ 沿候选行广播。};}
\newcommand{\TensorLegend}[2]{%
  \foreach \x/\col/\txt in {.55/Data/{输入 / 二值支持},5.1/Hidden/{隐藏激活},9.65/Weight/{可训练参数},14.2/Score/{分数 / 损失},18.75/Target/{位置 / 标签 / ID},24.1/Chosen/{选中 ID}}{%
    \path[fill=\col!65,draw=\col!45,rounded corners=.6pt]
      (\x,#1-.1) rectangle ++(.2,.2);
    \node[note,anchor=west] at (\x+.33,#1) {\txt};}
  \draw[flow] (.55,#1-.49) -- (1.5,#1-.49);
  \node[note,anchor=west] at (1.68,#1-.49) {实线：前向 / 数据流};
  \draw[grad] (10.3,#1-.49) -- (9.35,#1-.49);
  \node[note,anchor=west] at (10.55,#1-.49) {红色虚线：梯度，方向与前向相反（仅训练）};
  \node[note,anchor=west,text width=27cm,align=left] at (.55,#1-1.02) {#2};}

\hypersetup{pdfauthor={shichao},pdftitle={Metis training tensor flow},
  pdfsubject={Tensor dimensions, ranking loss and LoRA gradient direction},
  pdfkeywords={Metis, tensors, training, gradients, ScoreHead},pdfcreator={TeX}}
\begin{document}
\sffamily
\begin{tikzpicture}[x=1cm,y=1cm,text=Ink,font=\Body]
\path[use as bounding box] (0,0) rectangle (28,27.3);
\node[anchor=west,font=\fontsize{15}{18}\selectfont\bfseries]
  at (.25,26.75) {Metis 训练：候选分数、排序损失与反向梯度};
\node[anchor=west,note] at (.25,26.18)
  {$b\leq4$：候选微批；$N\leq8$：每个 query 候选数；$J\leq N$：已标注数。};
\node[anchor=east,note] at (27.75,26.18)
  {$L\leq2048,\quad D=1024,\quad h=256$；query batch $=1$};

\BackbonePanel{22.5}{\texttt{train()}；保留激活计算图}
\HeadPanel{16.3}

% Score/loss panel: the only branches are local, straight diagonals.
\Band{.1}{7.0}{27.9}{13.2}{Score}
\node[anchor=west,heading,text=Score!85!Ink] at (.4,12.76)
  {3. 汇总当前 query 的微批，计算排序损失（FP32）};
\node[anchor=east,note] at (27.5,12.76) {未标注候选不补零标签};
\Tensor{part}{1.1}{10.8}{.4}{1.05}{3}{1}{Score}{\boldsymbol{s}_{\mathrm{chunk}}}
\Dim{1.1}{9.48}{[b]}
\node[block,minimum width=1.8cm,minimum height=1cm,font=\Small]
  (cat) at (3.55,10.8) {\textbf{cat}\\[4pt]全部微批};
\Tensor{scores}{6.05}{10.8}{.4}{2.1}{6}{1}{Score}{\boldsymbol{s}}
\Dim{6.05}{8.98}{[N]}
\node[block,minimum width=2.5cm,minimum height=1.1cm,font=\Small]
  (known) at (8.65,10.8) {取已标注位置\\[4pt]$t=s[\mathrm{known}]$};
\Tensor{judged}{11.2}{10.8}{.45}{1.05}{3}{1}{Score}{t}
\Dim{11.2}{9.48}{[J]}
\node[note] at (8.65,9.43) {$\mathrm{known}=(0,2,4)$};
\Tensor{grades}{11.2}{8.6}{.45}{1.05}{3}{1}{Target}{g}
\Dim{11.2}{7.35}{[J]}
\node[note] at (8.35,8.5) {独立的监督标签\\不进入模型输入};
\draw[flow] (part-east) -- (cat.west);
\draw[flow] (cat.east) -- (scores-west);
\draw[flow] (scores-east) -- (known.west);
\draw[flow] (known.east) -- (judged-west);

\node[block,minimum width=1.2cm,minimum height=1.1cm,font=\Small]
  (inputs) at (13.0,10.0) {$t,g$};
\draw[flow] (judged-east) -- ($(inputs.west)+(0,.3)$);
\draw[flow] (grades-east) -- ($(inputs.west)+(0,-.3)$);
\node[block,minimum width=8.0cm,minimum height=1.05cm,font=\Small]
  (bce) at (18.15,11.65)
  {$\ell_B=\operatorname{mean}\,\operatorname{BCEWithLogits}(t,\mathbf{1}[g>0])$};
\node[note] at (18.15,12.36) {二元相关性，$[J]\to[\,]$};
\node[block,minimum width=8.0cm,minimum height=3.35cm]
  (pair) at (18.15,8.89) {};
\node[note] at (18.15,10.21) {严格等级偏好，$[\#\mathrm{pairs}]\to[\,]$};

% Exact illustrative supports: g=(2,1,0), t=(1.2,.4,-.7).
\foreach \r in {0,1,2}{\foreach \c in {0,1,2}{%
  \ifnum\r<\c\def\support{70}\else\def\support{0}\fi
  \path[draw=Data!20,fill=Data!\support,line width=.2pt]
    ({14.6+\c*.34},{8.55+(2-\r)*.34}) rectangle ++(.30,.30);
  \ifnum\r=\c\def\shade{0}\def\sign{0}
  \else\ifnum\r<\c\def\shade{65}\def\sign{+}
  \else\def\shade{20}\def\sign{-}\fi\fi
  \path[draw=Score!20,fill=Score!\shade,line width=.2pt]
    ({16.55+\c*.34},{8.55+(2-\r)*.34}) rectangle ++(.30,.30);
  \node[font=\fontsize{7}{8}\selectfont] at ({16.7+\c*.34},{8.70+(2-\r)*.34}) {$\sign$};
}}
\node[font=\Small] at (15.09,8.27) {$C$};
\Dim{15.09}{7.88}{[J,J]\;\mathrm{bool}}
\node[font=\Small] at (17.04,8.27) {$\Delta$};
\Dim{17.04}{7.88}{[J,J]}
\node[note,align=left] at (19.9,9.03)
  {$C_{ij}=\mathbf{1}[g_i>g_j]$\\[5pt]$\Delta_{ij}=t_i-t_j$};
\node[font=\Small] at (18.15,7.5)
  {$\ell_P=\operatorname{mean}\,\operatorname{softplus}(-\Delta[C])$};
\draw[flow] ($(inputs.east)+(0,.3)$) -- (bce.west);
\draw[flow] ($(inputs.east)+(0,-.3)$) -- (pair.west);

\node[block,draw=Score!45,fill=Score!8,minimum width=3.7cm,
  minimum height=2.1cm,font=\Body] (loss) at (25.45,10.1)
  {$\mathcal L_q$\\[5pt]$0.1\ell_B+\ell_P$\\[5pt]$[\,]$ 标量};
\draw[flow] (bce.east) -- ($(loss.west)+(0,.45)$);
\draw[flow] (pair.east) -- ($(loss.west)+(0,-.45)$);
\node[note] at (25.45,8.22)
  {$\mathcal L=\operatorname{mean}_q\mathcal L_q$\\[5pt]无偏好对时 $\ell_P=0$};
\node[note,anchor=west] at (.4,7.66)
  {格子示例：$t=(1.2,0.4,-0.7)$，$g=(2,1,0)$。\\
   $C$ 白格为 $0$；$\Delta$ 内的 $+/-/0$ 仅说明示例符号。};

% Dedicated backward lane; arrows go right to left, without page-spanning loops.
\Band{.1}{3.0}{27.9}{6.65}{Score}
\node[anchor=west,heading,text=Score!85!Ink] at (.4,6.22)
  {4. 反向传播：按上方同名张量对应，从右向左};
\node[block,draw=Score!40,minimum width=2.6cm,minimum height=1.15cm,font=\Small]
  (seed) at (25.45,4.8) {$\mathcal L$\\[4pt]标量损失};
\node[block,draw=Score!40,minimum width=3.35cm,minimum height=1.15cm,font=\Small]
  (gradscore) at (20.9,4.8)
  {$\partial\mathcal L/\partial\boldsymbol{s}\;[N]$\\[4pt]cat 按微批切回 $[b]$};
\node[block,draw=Weight!45,fill=Weight!7,minimum width=3.5cm,
  minimum height=1.15cm,font=\Small] (gradhead) at (16.35,4.8)
  {Head 反向\\[4pt]$\nabla_{\phi}\mathcal L$};
\node[block,draw=Hidden!40,minimum width=3cm,minimum height=1.15cm,font=\Small]
  (gradZ) at (12.05,4.8) {$\partial\mathcal L/\partial Z$\\[4pt]$[b,D]$};
\node[block,draw=Hidden!40,minimum width=3cm,minimum height=1.15cm,font=\Small]
  (gradH) at (7.95,4.8) {$\partial\mathcal L/\partial H$\\[4pt]$[b,L,D]$};
\node[block,draw=Weight!45,fill=Weight!7,minimum width=3.1cm,
  minimum height=1.15cm,font=\Small] (gradlora) at (3.2,4.8)
  {Decoder / LoRA\\[4pt]$\nabla_{\Delta\theta}\mathcal L$};
\draw[grad] (seed.west) -- (gradscore.east);
\draw[grad] (gradscore.west) -- (gradhead.east);
\draw[grad] (gradhead.west) -- (gradZ.east);
\draw[grad] (gradZ.west) -- (gradH.east);
\draw[grad] (gradH.west) -- (gradlora.east);
\node[note,anchor=west] at (.4,3.47)
  {更新 $\Delta\theta$ 与 $\phi=(W_1,b_1,W_2,b_2)$，参数梯度同形；$\theta_0$ 冻结但激活梯度继续传播。$X,M,P,g,C,\mathrm{known}$ 不求梯度。};

\TensorLegend{1.85}{FP32 主参数与更新；BF16 autocast 用于支持的计算，激活 dtype 不统一假定为 BF16。分数与损失为 FP32。跨行同名为同一张量的引用，未 detach；cat 保留顺序和计算图。色格数量、深浅仅作示意。}
\node[note] at (14,.18) {Metis 训练张量与梯度 · 首轮 pairs / SDPA 配方 · shichao};
\end{tikzpicture}
\end{document}
